% Memo to DWTS Producers (1-2 pages as required by problem statement)

\begin{letter}{Memo to DWTS Producers}
\label{sec:memo}

\begin{center}
\rule{0.9\textwidth}{1.5pt}\\[0.2cm]
{\Large\bfseries MEMORANDUM}\\[0.2cm]
\rule{0.9\textwidth}{1.5pt}
\end{center}

\vspace{0.2cm}

\noindent
\begin{tabular}{@{}ll}
\textbf{TO:} & Executive Producers, \textit{Dancing with the Stars} \\
\textbf{FROM:} & Team \#2603956, MCM Data Analysis Group \\
\textbf{DATE:} & February 2026 \\
\textbf{RE:} & Recommendations on Combining Judge Scores and Fan Votes \\
\end{tabular}

\vspace{0.3cm}

\subsection*{Executive Summary}
We analyzed 34 seasons of DWTS data to address: \textbf{how should judge scores and fan votes be combined to produce fair and engaging outcomes?} Our analysis reconstructs hidden fan vote shares, compares aggregation rules, and examines controversial cases (Bobby Bones in Season~27). We recommend a hybrid policy balancing audience engagement with technical credibility.

\subsection*{Key Findings}
\begin{enumerate}
\setlength{\itemsep}{0.3ex}
    \item \textbf{Fan votes can be reliably inferred.} Using elimination-consistency constraints, we reconstructed weekly fan vote shares that reproduce 91\% of historical eliminations.
    \item \textbf{The percentage rule better reflects audience preference.} When the two rules disagree, percentage more often retains contestants with higher fan support.
    \item \textbf{Judge-save reduces controversy but shifts power.} Introducing judge-selected elimination prevents outcomes like Bobby Bones' victory, but transfers significant influence to judges.
    \item \textbf{Judges and fans weight attributes differently.} Judges emphasize progress and partner skill; fans respond more to popularity and celebrity identity.
\end{enumerate}

\subsection*{Recommendations}
\begin{enumerate}
\setlength{\itemsep}{0.3ex}
    \item[\textbf{R1.}] \textbf{Adopt a hybrid aggregation policy.}
    Use \textit{percentage-based} aggregation in early weeks (Weeks~1--6) to maximize audience engagement. Switch to \textit{judge-selected bottom-three elimination} at mid-season (Week~7 onward) to ensure technical quality influences final placements.

    \item[\textbf{R2.}] \textbf{Monitor fan--judge divergence in real time.} Large divergences (top-3 fan support but bottom-3 judge ranking) are early warning signs of potential controversy.

    \item[\textbf{R3.}] \textbf{Use transparency strategically.} Highlight ``divergence weeks'' to enhance narrative tension while maintaining trust.
\end{enumerate}

\subsection*{Why This Approach Works}
\begin{itemize}
\setlength{\itemsep}{0.2ex}
    \item \textbf{Audience satisfaction:} Fans see their votes matter early; late outcomes reflect both popularity and skill.
    \item \textbf{Technical credibility:} Avoids outcomes where low-scoring contestants win (Season~27).
    \item \textbf{Production simplicity:} Fits existing scoring infrastructure---no new voting technology required.
    \item \textbf{Narrative value:} Divergence between fans and judges creates natural storylines that can be highlighted rather than hidden.
\end{itemize}

\subsection*{Implementation Considerations}
The hybrid policy can be implemented immediately within the existing production framework. Early-season percentage aggregation requires no changes to current voting infrastructure. The mid-season transition to judge-selected elimination mirrors the existing dance-off format, where judges already make subjective decisions. We recommend announcing the rule structure at the season premiere to set viewer expectations and enhance transparency. Real-time divergence monitoring can be achieved using the same data pipelines that currently track votes and scores.

\vspace{0.2cm}
\noindent\textit{Respectfully submitted,}\\
\textbf{Team \#2603956}

\end{letter}
